% =============================================================================
% SunGauss: Physically-Based Outdoor Scene Relighting with Gaussian Splatting
% DRAFT — February 2026
% =============================================================================
\documentclass[conference]{IEEEtran}
% \documentclass{article} % alternative: switch to your venue class

% ── packages ────────────────────────────────────────────────────────────────
\usepackage[utf8]{inputenc}
\usepackage{amsmath,amssymb,amsfonts}
\usepackage{graphicx}
\usepackage{booktabs}
\usepackage{hyperref}
\usepackage{xcolor}
\usepackage{algorithm}
\usepackage{algorithmic}
\usepackage{multirow}
\usepackage{subcaption}
\usepackage{siunitx}
\usepackage{dblfloatfix}

\setcounter{dbltopnumber}{2}
\renewcommand{\dbltopfraction}{0.95}
\renewcommand{\dblfloatpagefraction}{0.85}
\renewcommand{\textfraction}{0.05}

% draft helpers
\newcommand{\todo}[1]{{\color{red}\textbf{[TODO:~#1]}}}
\newcommand{\ours}{\textsc{SunGauss}}

% ── title ───────────────────────────────────────────────────────────────────
\title{\ours{}: Physically-Based Sun Relighting for\\In-the-Wild Outdoor Scenes via Gaussian Splatting}

\author{
    \IEEEauthorblockN{Alp Bora Orgun, Marco Volino, Adrian Hilton, Jean-Yves Guillemaut}
    \IEEEauthorblockA{CVSSP, University of Surrey\\
    \texttt{ao00821@surrey.ac.uk, m.volino@surrey.ac.uk, a.hilton@surrey.ac.uk, j.guillemaut@surrey.ac.uk}}
}

\begin{document}
\maketitle

% ============================================================================
% ABSTRACT
% ============================================================================
\begin{abstract}
    We present \ours{}, a method for relighting outdoor scenes captured in uncontrolled conditions under natural sunlight.
    Built upon 2D Gaussian Splatting (2DGS) with appearance embeddings, we replace the per-image spherical-harmonic environment maps of prior work with an explicit physically-motivated sun lighting model that decomposes illumination into a direct solar term, scene-global ambient light parameterised as spherical harmonics of the sun direction, and a shared global sky component.
    Sun intensity, colour correction, and ambient are all represented as low-order SH functions of the sun direction, reducing lighting parameters to $O(1)$ and enabling evaluation at novel sun positions unseen during training.
    Our explicit decomposition enables several capabilities absent from the baseline:
    (i)~geometry-based shadow computation via differentiable shadow mapping from an orthographic sun camera,
    and (ii)~realistic day-cycle relighting by simply sweeping the sun direction through a solar trajectory.
    We further introduce learnable per-camera sun direction calibration and a per-Gaussian sky distinction parameter supervised by binary sky masks.
    On the NeRF-OSR benchmark, \ours{} outperforms NeRF-OSR on all three evaluated scenes and is competitive with LumiGauss, achieving higher PSNR on Landwehrplatz and Staatstheater (18.91 and 17.98 dB versus 18.01 and 17.02), lower MSE on two scenes, and improved SSIM on Ludwigskirche and Staatstheater (0.730 and 0.753 versus 0.700 and 0.729), while enabling explicit geometry-based shadows and relighting at novel sun directions.
\end{abstract}

% ============================================================================
\section{Introduction}
\label{sec:intro}
% ============================================================================

Relighting outdoor scenes from uncontrolled photo collections is a long-standing challenge in computer vision and graphics.
Unlike studio environments with known illumination, in-the-wild captures exhibit varying sun position, cloud cover, and atmospheric conditions across images.
Recent neural scene representations---particularly Neural Radiance Fields (NeRFs)~\cite{mildenhall2020nerf} and Gaussian Splatting~\cite{kerbl20233dgs}---have enabled photorealistic novel-view synthesis, but most methods bake illumination into the scene representation, preventing post-capture relighting.

Prior works such as LumiGauss~\cite{joaxkal2024lumigauss} addresses this for outdoor scenes by combining 2D Gaussian Splatting (2DGS)~\cite{huang20242dgs} with per-image appearance embeddings and spherical-harmonic (SH) environment maps.
Each image is assigned a learnable SH environment that modulates Gaussian colors, enabling factorization of geometry, albedo, and illumination.
However, the per-image SH representation has important limitations for outdoor relighting:
\begin{itemize}
    \item SH environments are unconstrained and per-image---they do not share structure across images, making it impossible to interpolate to novel, unseen illumination conditions.
    \item Shadows are encoded only implicitly through per-Gaussian SH feature modulation, which cannot produce sharp cast shadows from geometry.
    \item There is no explicit notion of the sun as a directional source, precluding physically-grounded relighting such as time-of-day simulation.
\end{itemize}

We address these limitations with \ours{}, which makes the following contributions:

\begin{enumerate}
    \item \textbf{Explicit directional sun model with scene-global SH lighting} (\S\ref{sec:sun_model}): A physically-motivated lighting decomposition separating direct sunlight (with Lambert shading and an atmospheric colour prior), ambient light, and a shared global sky SH.  Sun intensity, colour correction, and ambient are represented as low-order SH functions of the sun direction, reducing lighting parameters to $O(1)$, enforcing physical consistency across views, and enabling evaluation at novel sun positions.

    \item \textbf{Geometry-based shadow computation} (\S\ref{sec:shadows}): Shadow mapping via orthographic sun projection applied to the direct lighting term.

    \item \textbf{Learnable sun direction calibration} (\S\ref{sec:sun_cal}): Per-camera learnable offsets that jointly refine noisy input sun directions during training.

    \item \textbf{Learnable sky-scene distinction} (\S\ref{sec:sky_mask}): A per-Gaussian continuous parameter supervised by sky masks that controls both shadow casting and shadow receiving.
\end{enumerate}

% ============================================================================
\section{Related Work}
\label{sec:related}
% ============================================================================

\paragraph{Neural Scene Relighting.}
\todo{}
NeRF-OSR~\cite{rudnev2022nerfOSR} decomposes outdoor NeRF scenes into albedo and per-image SH lighting, enabling appearance transfer between views.
NeRF-W~\cite{martinbrualla2021nerfw} handles appearance variation with per-image latent codes.
More recently, methods based on 3D Gaussian Splatting~\cite{kerbl20233dgs} have achieved real-time rendering while maintaining relighting capability.

\paragraph{Outdoor Illumination Modeling.}
\todo{}Modeling outdoor illumination typically involves representing the sun as a directional source and the sky as a hemisphere of diffuse light~\cite{lalonde2012estimating,hold-geoffroy2019deep}.
The color of sunlight varies with elevation due to atmospheric Rayleigh scattering, producing the characteristic warm tones at sunrise and sunset~\cite{preetham1999practical,nishita1993display}.
Our sun color prior draws on this atmospheric model.

\paragraph{Shadow Computation.}
\todo{}Shadow mapping~\cite{williams1978casting} renders a depth map from the light's viewpoint and compares fragment depths to determine occlusion.
In the Gaussian splatting setting, the unstructured point-based representation requires adapting shadow mapping to work with splatted depth rather than mesh rasterization.

\paragraph{2D Gaussian Splatting.}
\todo{}2DGS~\cite{huang20242dgs} represents scene geometry as oriented planar splats (surfels), providing well-defined surface normals from the splat orientation.
LumiGauss~\cite{joaxkal2024lumigauss} builds on 2DGS, leveraging the surfel normals for hemisphere-based SH evaluation.
We inherit this representation and additionally leverage the normals for Lambert shading (N$\cdot$L).

% ============================================================================
\section{Method}
\label{sec:method}
% ============================================================================

Given a set of posed images $\{I_i\}_{i=1}^{N}$ of an outdoor scene with known per-image sun directions $\{\mathbf{l}_i\}$ and optional binary sky masks $\{M_i^{\text{sky}}\}$, our goal is to learn a Gaussian scene representation that decomposes geometry, material (albedo), and illumination, enabling relighting under novel sun conditions.

We build on the 2DGS surfel representation~\cite{huang20242dgs} with per-image appearance embeddings from LumiGauss~\cite{joaxkal2024lumigauss}.
Our method replaces the per-image SH environment with an explicit sun lighting model, adds geometry-based shadow computation, and introduces several training improvements.
Figure~\ref{fig:pipeline} provides an overview.

\begin{figure*}[t]
    \centering
    \todo{Pipeline figure: show Gaussians $\rightarrow$ normals $\rightarrow$ N$\cdot$L $\rightarrow$ shadow map $\rightarrow$ composited image}
    \caption{Overview of the \ours{} pipeline. Per-Gaussian normals from the surfel orientation are used for Lambert shading with the sun direction. A shadow map rendered from an orthographic sun camera provides per-Gaussian shadow masks. The shadowed direct light, ambient, and sky SH components are composed and multiplied by learned albedo to produce the final image.}
    \label{fig:pipeline}
\end{figure*}

% ── 3.1 Sun Model ──────────────────────────────────────────────────────────
\subsection{Explicit Directional Sun Model}
\label{sec:sun_model}

We decompose the illumination reaching each Gaussian into three physically-motivated terms:

\begin{equation}
    L(\mathbf{x}, \mathbf{n}) = \underbrace{C_{\text{sun}}(\theta) \cdot I(\mathbf{l}_i) \cdot \mathbf{k}(\mathbf{l}_i) \cdot \max(0, \mathbf{n} \cdot \mathbf{l}_i)}_{\text{direct sunlight}} + \underbrace{\mathbf{c}^{\text{amb}}(\mathbf{l}_i)}_{\text{ambient}} + \underbrace{E_{\text{sky}}(\mathbf{n})}_{\text{global sky}}
    \label{eq:lighting}
\end{equation}

\noindent where $\mathbf{n}$ is the surface normal (derived from the surfel quaternion and scale, \S\ref{sec:normals}), $\mathbf{l}_i$ is the sun direction for image $i$, and:

\paragraph{Sun color prior $C_{\text{sun}}(\theta)$.}
We model the spectral shift of sunlight with elevation $\theta$ using an atmospheric Rayleigh scattering approximation.
The raw transmittance per RGB channel is:
\begin{equation}
    \hat{T}_c = \exp\!\left(-\beta_R \cdot \text{AM}(\theta) \cdot \left(\frac{\lambda_c}{650}\right)^{-4} \cdot s\right)
    \label{eq:sun_color}
\end{equation}
where $\text{AM}(\theta) = 1/\sin(\theta)$ is the air mass (clamped to $10 + 2(5 - \theta)$ for $\theta < 5^\circ$ to avoid singularity at the horizon), $\lambda_c \in \{650, 550, 450\}$\,nm for R, G, B, $\beta_R = 0.0596$ is the Rayleigh extinction coefficient, and $s = 0.15$ is a scattering scale.
To ensure the prior yields neutral white light at zenith, we normalise by the zenith ($\theta = 90^\circ$) transmittance:
\begin{equation}
    C_{\text{sun},c}(\theta) = \operatorname{clamp}\!\left(\frac{\hat{T}_c(\theta)}{\hat{T}_c(90^\circ)},\; 0.1,\; 1.5\right)
\end{equation}
This produces physically plausible warm tones at low elevations and neutral white at zenith.

\paragraph{Scene-global SH lighting parameters.}
Rather than storing per-image scalars---which cannot generalise to unseen sun positions---we represent the three lighting quantities as low-order spherical-harmonic functions of the \emph{sun direction} $\mathbf{l}$:
\begin{align}
    I(\mathbf{l})                       & = \sum_{l=0}^{d}\sum_{m=-l}^{l} a_{lm}\, Y_l^m(\hat{\mathbf{l}}), \label{eq:scene_sh_int}          \\
    \mathbf{k}(\mathbf{l})              & = \sum_{l=0}^{d}\sum_{m=-l}^{l} \mathbf{b}_{lm}\, Y_l^m(\hat{\mathbf{l}}), \label{eq:scene_sh_col} \\
    \mathbf{c}^{\text{amb}}(\mathbf{l}) & = \sum_{l=0}^{d}\sum_{m=-l}^{l} \mathbf{d}_{lm}\, Y_l^m(\hat{\mathbf{l}}), \label{eq:scene_sh_amb}
\end{align}
where $d = 2$ (9 SH coefficients each) by default, giving:
\begin{itemize}
    \item $I(\mathbf{l}) \in \mathbb{R}^+$: sun intensity multiplier (DC initialized to $4.0 / c_0$, $c_0 = 0.282095$)
    \item $\mathbf{k}(\mathbf{l}) \in [0.5, 2.0]^3$: colour correction on $C_{\text{sun}}$, accommodating cloud cover and atmospheric variation (DC initialized to $1.0 / c_0$)
    \item $\mathbf{c}^{\text{amb}}(\mathbf{l}) \in \mathbb{R}^{3}_+$: ambient/indirect illumination (DC initialized to $0.15 / c_0$)
\end{itemize}
Higher-order SH terms start at zero.
The effective sun colour is $C_{\text{sun}}(\theta) \cdot \mathbf{k}(\mathbf{l}) \cdot I(\mathbf{l})$.
This reduces the total lighting parameters to a fixed 39 scalars (9 for intensity, $3{\times}9$ for colour correction and ambient), enforces physical consistency across views with similar sun positions, and---crucially---enables evaluation at novel sun directions unseen during training.

\paragraph{Global sky SH.}
We represent the directional sky illumination as degree-1 spherical harmonics $E_{\text{sky}} \in \mathbb{R}^{3 \times 4}$, shared across all images.
This captures the dominant sky gradient (bright horizon, dark zenith) while remaining low-dimensional enough to constrain the optimization.
Evaluated at normal $\mathbf{n}$:
\begin{equation}
    E_{\text{sky}}(\mathbf{n}) = \sum_{l=0}^{1}\sum_{m=-l}^{l} c_{lm} \cdot Y_l^m(\mathbf{n})
    \label{eq:sky_sh}
\end{equation}

\paragraph{Final per-Gaussian color.}
The HDR intensity is gamma-corrected and multiplied by the learned albedo $\mathbf{a}$:
\begin{equation}
    \text{RGB} = \mathbf{a} \cdot \left(L(\mathbf{x}, \mathbf{n})\right)^{1/\gamma}, \quad \gamma = 2.2
    \label{eq:final_color}
\end{equation}

% ── 3.2 Normals ───────────────────────────────────────────────────────────
\subsection{Surface Normals from Surfel Geometry}
\label{sec:normals}

Each 2DGS surfel is parameterized by a quaternion $\mathbf{q}$ and anisotropic scales $\mathbf{s} \in \mathbb{R}^3$.
The surfel plane normal is the third column of the rotation matrix $R(\mathbf{q})$:
\begin{equation}
    \mathbf{n}_w = R(\mathbf{q})_{:,3}
\end{equation}
We apply a view-dependent sign flip to ensure normals face the camera:
\begin{equation}
    \mathbf{n} = \text{sign}(-\mathbf{n}_w \cdot \mathbf{p}_{\text{view}}) \cdot \mathbf{n}_w
\end{equation}
where $\mathbf{p}_{\text{view}}$ is the Gaussian position in camera coordinates.
This yields consistent camera-facing normals for Lambert shading.

% ── 3.3 Shadow Computation ─────────────────────────────────────────────────
\subsection{Shadow Computation}
\label{sec:shadows}

Shadows are applied exclusively to the direct sunlight term; ambient and sky illumination remain unaffected.
The shadowed intensity for Gaussian $j$ in image $i$ becomes:
\begin{equation}
    L_{\text{shad}} = C_{\text{sun}} \cdot I(\mathbf{l}_i) \cdot \mathbf{k}(\mathbf{l}_i) \cdot (\mathbf{n}\cdot\mathbf{l}_i) \cdot \sigma_j + \mathbf{c}^{\text{amb}}(\mathbf{l}_i) + E_{\text{sky}}(\mathbf{n})
    \label{eq:shadowed}
\end{equation}
where $\sigma_j \in [0, 1]$ is the shadow mask (1 = lit, 0 = shadowed).

We implement shadow computation using shadow mapping.

\subsubsection{Shadow Mapping}
\label{sec:shadow_map}

Our primary method adapts classical shadow mapping~\cite{williams1978casting} to Gaussian splatting.
We construct an orthographic sun camera positioned along $\mathbf{l}_i$ at distance $3\times r_{\text{scene}}$ from the scene center, with frustum half-width $1.5\times r_{\text{scene}}$, where $r_{\text{scene}}$ is the maximum distance from any Gaussian to the scene centroid.

We render this view using the surfel rasterizer to obtain a depth map $D_{\text{sun}} \in \mathbb{R}^{R \times R}$ and alpha map $\alpha_{\text{sun}}$.
Sky Gaussians are excluded by zeroing their opacity using the learned $w_j^{\text{sky}}$ parameter (\S\ref{sec:sky_mask}).

For each Gaussian $j$ with world position $\mathbf{x}_j$, we project into the sun camera's clip space, bilinearly sample the shadow depth map, and compare:
\begin{equation}
    \sigma_j = \begin{cases}
        0 & \text{if } d_j > D_{\text{sun}}(\mathbf{u}_j) + b \;\wedge\; \alpha_{\text{sun}}(\mathbf{u}_j) > 0.1 \\
        1 & \text{otherwise}
    \end{cases}
\end{equation}
where $d_j$ is the projected depth, $\mathbf{u}_j$ is the UV coordinate, and $b$ is a bias term to prevent self-shadowing.
Gaussians outside the shadow frustum default to lit.

% ── 3.3b Sun Direction Calibration ─────────────────────────────────────────
\subsection{Learnable Sun Direction Calibration}
\label{sec:sun_cal}

Sun directions from metadata (e.g.\ solar position databases or SfM-derived estimates) may be noisy or biased.
We optionally refine them during training by attaching a learnable offset $\Delta\mathbf{l}_i \in \mathbb{R}^3$ (initialized to zero) to each camera:
\begin{equation}
    \hat{\mathbf{l}}_i = \frac{\mathbf{l}_i + \Delta\mathbf{l}_i}{\|\mathbf{l}_i + \Delta\mathbf{l}_i\|}
    \label{eq:sun_cal}
\end{equation}
The refined direction $\hat{\mathbf{l}}_i$ replaces $\mathbf{l}_i$ in all downstream computations (Lambert shading, shadow mapping, scene-SH evaluation).
To prevent unconstrained drift, we apply an L2 regulariser:
\begin{equation}
    \mathcal{L}_{\text{sun-cal}} = \lambda_{\text{sc}} \|\Delta\mathbf{l}_i\|^2
    \label{eq:sun_cal_reg}
\end{equation}
Calibration is active only in the window $[T_{\text{sc,start}},\, T_{\text{sc,end}}]$ (default 10\,k--30\,k iterations), giving the geometry time to stabilise before the sun direction is adjusted, and freezing the directions before the final refinement iterations.

% ── 3.4 Learnable Sky Mask ──────────────────────────────────────────────────
\subsection{Learnable Sky-Scene Distinction}
\label{sec:sky_mask}

Each Gaussian carries a continuous parameter $w_j^{\text{sky}} \in [0, 1]$, implemented as a clamped learnable scalar.
This parameter serves a dual role:
\begin{itemize}
    \item \textbf{Shadow casting}: In the shadow map, each Gaussian's opacity is multiplied by $w_j^{\text{sky}}$, so sky Gaussians ($w_j^{\text{sky}} \approx 0$) become transparent to the sun.
    \item \textbf{Shadow receiving}: Sky Gaussians are forced to fully lit ($\sigma_j = 1$) regardless of the shadow computation, preventing the sky from being darkened by shadows.
\end{itemize}

\paragraph{Supervision.}
Given binary sky masks $M^{\text{sky}}$ (1 = non-sky, 0 = sky), we render each view using $w_j^{\text{sky}}$ as the per-Gaussian color and supervise with an L1 loss:
\begin{equation}
    \mathcal{L}_{\text{sky}} = \lambda_{\text{sky}} \left\| \text{Render}(w^{\text{sky}}) - M^{\text{sky}} \right\|_1
    \label{eq:sky_loss}
\end{equation}
This is differentiable through the rasterizer and jointly optimized with all other parameters.

\paragraph{Sky distance regularisation.}
To prevent sky Gaussians from drifting into the scene interior, we add a penalty pushing them away from the scene centre:
\begin{equation}
    \mathcal{L}_{\text{sky-dist}} = \lambda_{\text{sd}} \frac{1}{|\mathcal{S}|}\sum_{j \in \mathcal{S}} \left[\max(0,\; f \cdot r_{\text{scene}} - \|\mathbf{x}_j\|)\right]^2
    \label{eq:sky_dist}
\end{equation}
where $\mathcal{S} = \{j : w_j^{\text{sky}} < 0.5\}$ is the set of identified sky Gaussians, $r_{\text{scene}}$ is the cameras extent, and $f = 4$ (default) is the minimum-distance factor.

\begin{table*}[!tb]
    \centering
    \caption{Quantitative comparison on NeRF-OSR scenes.}
    \label{tab:nvs}
    \setlength{\tabcolsep}{3pt}
    \small
    \begin{tabular}{l cccc cccc cccc}
        \toprule
                                                      & \multicolumn{4}{c}{Landwehrplatz} & \multicolumn{4}{c}{Ludwigskirche} & \multicolumn{4}{c}{Staatstheater}                                                                                    \\
        \cmidrule(lr){2-5} \cmidrule(lr){6-9} \cmidrule(lr){10-13}
        Method                                        & PSNR\,$\uparrow$                  & MSE\,$\downarrow$                 & MAE\,$\downarrow$                 & SSIM\,$\uparrow$
                                                      & PSNR\,$\uparrow$                  & MSE\,$\downarrow$                 & MAE\,$\downarrow$                 & SSIM\,$\uparrow$
                                                      & PSNR\,$\uparrow$                  & MSE\,$\downarrow$                 & MAE\,$\downarrow$                 & SSIM\,$\uparrow$                                                                 \\
        \midrule
        Yu et al.$_{\text{u/s}}$~\cite{yu2021outdoor} & 15.17                             & 0.033                             & 0.133                             & 0.376            & 17.87 & 0.017 & 0.097 & 0.378 & 15.28 & 0.032 & 0.138 & 0.385 \\
        Philip et al.~\cite{philip2019multiview}      & 12.28                             & 0.062                             & 0.179                             & 0.319            & 16.63 & 0.023 & 0.113 & 0.367 & 12.34 & 0.065 & 0.200 & 0.272 \\
        NeRF-OSR~\cite{rudnev2022nerfOSR}             & 16.65                             & 0.024                             & 0.114                             & 0.501            & 18.72 & 0.014 & 0.090 & 0.468 & 15.43 & 0.029 & 0.133 & 0.517 \\
        NeRF-OSR*~\cite{rudnev2022nerfOSR}            & 15.66                             & --                                & --                                & --               & 19.34 & 0.012 & --    & --    & 16.35 & 0.027 & --    & --    \\
        SR-TensoRF~\cite{jin2023tensoir}              & 16.74                             & 0.024                             & 0.093                             & 0.653            & 17.30 & 0.021 & 0.096 & 0.542 & 15.43 & 0.030 & 0.111 & 0.632 \\
        FEGR~\cite{fegr2024}                          & 17.57                             & 0.018                             & --                                & --               & 21.53 & 0.007 & --    & --    & 17.00 & --    & --    & --    \\
        SOL-NeRF~\cite{solnerf2024}                   & 17.58                             & 0.028                             & --                                & 0.618            & 21.23 & 0.008 & --    & \textbf{0.749} & \textbf{18.18} & 0.019 & --    & 0.680 \\
        NeuSky~\cite{neusky2024}                      & 18.31                             & 0.016                             & --                                & --               & \textbf{22.50} & \textbf{0.005} & --    & --    & 16.66 & 0.023 & --    & --    \\
        LumiGauss~\cite{joaxkal2024lumigauss}         & 18.01                             & 0.017                             & 0.096                             & \textbf{0.778}  & 19.59 & 0.012 & \textbf{0.085} & 0.700 & 17.02 & 0.021 & 0.107 & 0.729 \\
        Ours                                          & \textbf{18.91}                   & \textbf{0.015}                   & \textbf{0.091}                   & 0.751            & 19.08 & 0.013 & 0.086 & 0.730 & 17.98 & \textbf{0.017} & \textbf{0.096} & \textbf{0.753} \\
        % \midrule
        % Yu et al.~\cite{yu2021outdoor}                          & 15.84                             & 0.028                             & 0.134                             & 0.392            & 18.71   & 0.014   & 0.088   & 0.400   & 15.43   & 0.031   & 0.136   & 0.363   \\
        % Philip et al.$_{\text{d/s}}$~\cite{philip2019multiview} & 12.85                             & 0.054                             & 0.169                             & 0.414            & 17.37   & 0.018   & 0.110   & 0.429   & 11.85   & 0.073   & 0.210   & 0.184   \\
        % NeRF-OSR$_{\text{d/s}}$~\cite{rudnev2022nerfOSR}        & 17.38                             & 0.021                             & 0.106                             & 0.576            & 19.86   & 0.011   & 0.080   & 0.626   & 15.83   & 0.026   & 0.133   & 0.556   \\
        % LumiGauss$_{\text{d/s}}$~\cite{joaxkal2024lumigauss}    & 18.40                             & 0.016                             & 0.094                             & 0.746            & 20.13   & 0.011   & --      & 0.727   & 17.24   & 0.020   & 0.106   & 0.715   \\
        % \midrule
        % LumiGauss$\,\dagger$                                    & 15.03                             & 0.034                             & 0.139                             & 0.58             & 19.34   & 0.015   & 0.094   & 0.693   & 16.09   & 0.025   & --      & 0.665   \\
        % LumiGauss$\,\ddagger$                                   & 17.59                             & 0.019                             & 0.100                             & 0.733            & 19.05   & 0.016   & 0.097   & 0.680   & 16.83   & 0.022   & 0.110   & 0.694   \\
        % LumiGauss $\setminus\,\ell_{0.1}$                       & 18.30                             & 0.016                             & 0.095                             & 0.744            & 20.15   & 0.010   & 0.080   & 0.734   & 17.25   & 0.020   & 0.106   & 0.712   \\
        % LumiGauss $\setminus\,\ell_{s}$                         & 17.35                             & 0.020                             & 0.100                             & 0.728            & 20.17   & 0.012   & 0.081   & 0.729   & 17.10   & 0.020   & 0.106   & 0.703   \\
        \bottomrule
    \end{tabular}
\end{table*}


% ============================================================================
\section{Training}
\label{sec:training}
% ============================================================================

\todo{}

% Training proceeds in three phases.
% Crucially, the shadowed rendering loss is the \emph{primary} objective from the very first iteration: if the model were trained unshadowed against ground-truth images that contain real shadows, the optimisation would bake shadow patterns into the albedo since no shadow mask is available to explain the dark regions.

% \paragraph{Phase 1: Early warmup (iterations 0--$T_w$, default $T_w{=}20{,}000$).}
% The shadowed photometric loss (Eq.~\ref{eq:shadowed}) is active with weight $\lambda^{\text{shad}} = 1.0$.
% A small unshadowed regulariser ($\lambda^{\text{unshad}} = 0.01$) encourages clean albedo without overriding the shadow signal.
% 2DGS surfel regularisation ($\mathcal{L}_{\text{normal}},\, \mathcal{L}_{\text{dist}}$) is enabled after iteration~6\,000.

% \paragraph{Phase 2: Lighting refinement ($T_w$--$T_s$, default $T_s{=}20{,}500$).}
% The unshadowed regulariser is turned off ($\lambda^{\text{unshad}} = 0$), and only the shadowed loss drives the optimisation.
% Shadow maps and lighting parameters are refined jointly.

% \paragraph{Phase 3: Full training (iterations $> T_s$).}
% The shadowed loss remains the sole image-level objective ($\lambda^{\text{shad}} = 1.0$).
% Sun direction calibration (\S\ref{sec:sun_cal}), if enabled, is active from iteration~10\,k to~30\,k.

% \paragraph{Loss function.}
% The total loss combines:
% \begin{equation}
%     \mathcal{L} = \lambda^{\text{shad}} \mathcal{L}_{\text{img}}^{\text{shad}} + \lambda^{\text{unshad}} \mathcal{L}_{\text{img}}^{\text{unshad}} + \mathcal{L}_{\text{sky}} + \mathcal{L}_{\text{sky-dist}} + \mathcal{L}_{\text{sun-cal}} + \mathcal{L}_{\text{sun-reg}} + \mathcal{L}_{\text{surfel}}
%     \label{eq:total_loss}
% \end{equation}
% where $\mathcal{L}_{\text{img}} = (1-\lambda_{\text{ssim}})\mathcal{L}_1 + \lambda_{\text{ssim}}(1-\text{SSIM})$ is the standard photometric loss, $\mathcal{L}_{\text{sky}}$ is Eq.~\ref{eq:sky_loss}, $\mathcal{L}_{\text{sky-dist}}$ is Eq.~\ref{eq:sky_dist}, $\mathcal{L}_{\text{sun-cal}}$ is Eq.~\ref{eq:sun_cal_reg}, and:
% \begin{itemize}
%     \item \textbf{Surfel regularisation} $\mathcal{L}_{\text{surfel}}$: normal consistency and depth distortion losses from 2DGS~\cite{huang20242dgs}.
%     \item \textbf{Sun regularisation} $\mathcal{L}_{\text{sun-reg}}$: (a) a penalty if the ambient-to-sun intensity ratio exceeds 0.3, (b) L2 on the colour correction SH DC term toward $1.0/c_0$ and L2 on higher-order SH terms to keep corrections smooth, (c) L2 on sky SH coefficients to prevent degenerate solutions, (d) L2 on higher-order intensity and ambient SH terms, and (e) a penalty when sky SH DC magnitude exceeds half the ambient level.
% \end{itemize}

% ============================================================================
\section{Day-Cycle Relighting}
\label{sec:relighting}
% ============================================================================

Our explicit sun decomposition enables relighting by simply changing the sun direction $\mathbf{l}$ and elevation $\theta$.
We provide a day-cycle renderer that:
\begin{enumerate}
    \item Computes a solar trajectory from geographic coordinates (latitude, day-of-year) using standard solar position equations~\cite{meeus1991astronomical}---declination $\delta = 23.45^\circ \sin(360^\circ(284+d)/365)$, hour angle, and azimuth/elevation conversion
    \item For each sun position, evaluates the sun color prior (Eq.~\ref{eq:sun_color}), computes N$\cdot$L shading with the new direction, renders a shadow map, and composites the final image
    \item Produces a video sequence showing realistic time-of-day illumination changes with moving shadows
\end{enumerate}

Since the sky SH, albedo, and geometry are shared across all images, only the direct term and shadow map change.
The intensity, colour correction, and ambient are evaluated as SH functions of the new sun direction (Eqs.~\ref{eq:scene_sh_int}--\ref{eq:scene_sh_amb}), providing physically consistent lighting at sun positions never seen during training.

% ============================================================================
\section{Experiments}
\label{sec:experiments}
% ============================================================================

\subsection{Datasets}

We evaluate on the NeRF-OSR~\cite{rudnev2022nerfOSR} benchmark, which provides outdoor scenes captured under varying illumination with COLMAP reconstructions.
We use the \texttt{lk2} and \texttt{st} scenes with associated sun direction metadata and sky masks.

\subsection{Implementation Details}

\todo{}

% We train for 40{,}000 iterations with warmup $T_w = 20{,}000$ and shadow start $T_s = 20{,}500$.
% Shadow maps use $512 \times 512$ resolution (default) with depth bias $b = 0.1$.
% Sun model parameters are optimised with a base environment learning rate of $2 \times 10^{-2}$: intensity and ambient use $2{\times}$ this rate, colour correction uses $0.5{\times}$, and sky SH uses $1{\times}$.
% Albedo learning rate is $2.5 \times 10^{-3}$.
% When sun direction calibration is enabled, the per-camera deltas use learning rate $0.1$ with L2 weight $\lambda_{\text{sc}} = 10^{-4}$, active from iteration~10\,k to~30\,k.

\subsection{Baselines}

\begin{itemize}
    \item \textbf{LumiGauss}~\cite{joaxkal2024lumigauss}: Per-image SH environment maps with 2DGS. No explicit sun model or shadow computation.
    \item \textbf{NeRF-OSR}~\cite{rudnev2022nerfOSR}: NeRF-based outdoor relighting with SH decomposition.
    \item \todo{Additional baselines as needed.}
\end{itemize}

\subsection{Novel-View Synthesis}


\subsection{Relighting Quality}

\todo{Evaluate relighting with held-out illumination conditions. Show day-cycle results qualitatively. Compare shadow quality against baseline (implicit SH shadows).}

\subsection{Ablation Studies}

\todo{Add ablation studies.}

% ============================================================================
\section{Limitations and Future Work}
\label{sec:limitations}
% ============================================================================

\paragraph{Single directional light.}
Our model assumes a single dominant sun direction.
Overcast conditions with diffuse illumination are handled only through the ambient term, which may be insufficient for accurate relighting under heavily clouded skies or scenes with multiple separate major light sources.

\paragraph{Shadow map resolution.}
The discrete shadow map introduces aliasing at object boundaries.
Higher resolutions mitigate this at increased memory cost.
Percentage-closer filtering or variance shadow maps could improve edge quality.

\paragraph{Static geometry.}
Like the underlying Gaussian splatting, our method assumes a static scene.
Transient objects (pedestrians, vehicles) present in the training images are masked out.

\paragraph{Normal quality.}
Lambert shading quality depends on the accuracy of surfel normals.
In regions with poor multi-view coverage, normals may be noisy, leading to incorrect N$\cdot$L values and shading artifacts.

\paragraph{Future directions.}
Promising extensions include learnable shadow bias, cascaded shadow maps for multi-scale scenes, differentiable shadow computation for end-to-end optimization, and integration with sky models for full environment relighting including cloud rendering.

% ============================================================================
\section{Conclusion}
\label{sec:conclusion}
% ============================================================================

We presented \ours{}, a method for physically-based outdoor scene relighting using Gaussian splatting.
By replacing unconstrained per-image SH environments with an explicit directional sun model, we enable geometry-based shadow computation, atmospheric colour priors, and realistic day-cycle relighting---capabilities that are fundamentally out of reach for per-image latent representations.
The scene-global SH parameterisation of intensity, colour correction, and ambient further enables evaluation at novel sun directions unseen during training.
Learnable sun direction calibration and the learnable sky distinction further improve reconstruction quality.
\todo{Summary of quantitative improvements.}

% ============================================================================
% REFERENCES
% ============================================================================
{\small
    \bibliographystyle{IEEEtran}
    \begin{thebibliography}{99}

        \bibitem{mildenhall2020nerf}
        B.~Mildenhall, P.~P.~Srinivasan, M.~Tancik, J.~T.~Barron, R.~Ramamoorthi, and R.~Ng,
        ``NeRF: Representing scenes as neural radiance fields for view synthesis,''
        in \emph{ECCV}, 2020.

        \bibitem{kerbl20233dgs}
        B.~Kerbl, G.~Kopanas, T.~Leimk{\"u}hler, and G.~Drettakis,
        ``3D Gaussian Splatting for real-time radiance field rendering,''
        \emph{ACM Trans. Graphics}, vol.~42, no.~4, 2023.

        \bibitem{huang20242dgs}
        B.~Huang, Z.~Yu, A.~Chen, A.~Geiger, and S.~Gao,
        ``2D Gaussian Splatting for geometrically accurate radiance fields,''
        in \emph{SIGGRAPH}, 2024.

        \bibitem{joaxkal2024lumigauss}
        J.~Kaleta,
        ``LumiGauss: Relightable Gaussian Splatting in the Wild,''
        in \emph{WACV}, 2025.

        \bibitem{rudnev2022nerfOSR}
        V.~Rudnev, M.~Elgharib, W.~Smith, L.~Liu, V.~Golyanik, and C.~Theobalt,
        ``NeRF for outdoor scene relighting,''
        in \emph{ECCV}, 2022.

        \bibitem{martinbrualla2021nerfw}
        R.~Martin-Brualla, N.~Radwan, M.~S.~M.~Sajjadi, J.~T.~Barron, A.~Dosovitskiy, and D.~Duckworth,
        ``NeRF in the Wild: Neural radiance fields for unconstrained photo collections,''
        in \emph{CVPR}, 2021.

        \bibitem{williams1978casting}
        L.~Williams,
        ``Casting curved shadows on curved surfaces,''
        \emph{Computer Graphics (SIGGRAPH)}, vol.~12, no.~3, pp.~270--274, 1978.

        \bibitem{preetham1999practical}
        A.~J.~Preetham, P.~Shirley, and B.~Smits,
        ``A practical analytic model for daylight,''
        in \emph{SIGGRAPH}, 1999.

        \bibitem{nishita1993display}
        T.~Nishita, T.~Sirai, K.~Tadamura, and E.~Nakamae,
        ``Display of the earth taking into account atmospheric scattering,''
        in \emph{SIGGRAPH}, 1993.

        \bibitem{lalonde2012estimating}
        J.-F.~Lalonde, A.~A.~Efros, and S.~G.~Narasimhan,
        ``Estimating the natural illumination conditions from a single outdoor image,''
        \emph{IJCV}, vol.~98, no.~2, pp.~123--145, 2012.

        \bibitem{hold-geoffroy2019deep}
        Y.~Hold-Geoffroy, A.~Athawale, and J.-F.~Lalonde,
        ``Deep sky modeling for single image outdoor lighting estimation,''
        in \emph{CVPR}, 2019.

        \bibitem{meeus1991astronomical}
        J.~Meeus,
        \emph{Astronomical Algorithms}.
        Willmann-Bell, 1991.

        \bibitem{yu2021outdoor}
        \todo{Yu et al. -- outdoor relighting reference.}

        \bibitem{philip2019multiview}
        \todo{Philip et al. -- multi-view relighting reference.}

        \bibitem{jin2023tensoir}
        \todo{SR-TensoRF reference.}

        \bibitem{fegr2024}
        \todo{FEGR reference.}

        \bibitem{solnerf2024}
        \todo{SOL-NeRF reference.}

        \bibitem{neusky2024}
        \todo{NeuSky reference.}

    \end{thebibliography}
}

\end{document}
